\chapter{Fluxes, Tadpoles and \texorpdfstring{$\chi/24$}{chi/24}}\label{ch:tadpole}

On a compact space a charge cannot be alone. The field lines that leave it must end
somewhere, and if the space has no boundary and no infinity they can only end on other
charges of the opposite sign. This elementary remark about Gauss's law becomes, in string
theory, one of the sharpest consistency conditions a compactification has to pass: the
\emph{tadpole condition}. It is sharp because it is an equation between integers. For
M-theory or F-theory on a Calabi--Yau fourfold $X_8$ it reads
\begin{equation}\label{eq:tadpole:master}
 \frac{\chi(X_8)}{24}\;=\;n\;+\;\frac12\int_{X_8}\frac{G}{2\pi}\wedge\frac{G}{2\pi},
\end{equation}
where $\chi$ is the Euler number, $n$ counts space-filling branes and $G$ is the four-form
flux. This chapter answers four questions about \eqref{eq:tadpole:master}. Why is there such
an equation at all? Why does the curvature of the internal space enter through the number
$\chi/24$, and when is that number a positive integer? Why is the flux term an integer? And
what does the equation look like in the perturbative type IIB description of our main example,
the orientifold $\KT/\Z_2$ of \cref{ch:k3t2}, where the branes are D7- and D3-branes and the
negative charges sit on orientifold planes? For the product of two K3 surfaces every number
can be computed by hand: $\chi(K3\times K3)=576$, the budget is $576/24=24$, and the flux term
is an integer because the K3 lattice of \cref{ch:k3lattice} is even. These last statements
are arithmetic and lattice theory, and they are what the Lean library checks.

We use the following input. From \cref{ch:k3}: a K3 surface has Betti numbers $(1,0,22,0,1)$,
Hodge numbers $h^{0,0}=h^{2,0}=h^{0,2}=h^{2,2}=1$, $h^{1,1}=20$, and Euler number $24$. From
\cref{ch:k3lattice,ch:lattices}: $H^2(K3;\Z)$ with the cup product is the even unimodular
lattice $3U\oplus2E_8(-1)$ of signature $(3,19)$. From \cref{ch:dbranes}: a D$p$-brane is a
source of charge $\mu_p$ for the Ramond--Ramond potential $C_{p+1}$.

\section{Gauss's law on a compact space}\label{sec:tadpole:gauss}

\subsection{Electrodynamics first}
Let $M$ be a compact manifold without boundary, of dimension $d$, and let a gauge field with
field strength $F=\dd A$ be sourced by a charge density $\rho$. Gauss's law is the statement
\begin{equation}\label{eq:tadpole:gausslaw}
 \dd{\star F}\;=\;\rho\,\mathrm{vol}_M ,
\end{equation}
where $\star$ is the Hodge star of $M$ and $\mathrm{vol}_M$ its volume form. Integrate both
sides over $M$. The left-hand side is the integral of an exact form over a manifold without
boundary, which vanishes by Stokes's theorem. Hence
\begin{equation}\label{eq:tadpole:zero}
 Q_{\text{total}}=\int_M\rho\,\mathrm{vol}_M=0 .
\end{equation}
On $\R^3$ the same integration over a ball gives the flux of the electric field through the
boundary sphere, which can be anything; the difference is entirely the absence of a boundary.
Two features of this one-line argument will survive every generalization below. First, the
step uses only that $\star F$ is a \emph{globally defined} form: if it were defined only
patch by patch the integral of its differential need not vanish. Second, the conclusion is
insensitive to the metric; \eqref{eq:tadpole:zero} is a topological constraint on the sources.

\begin{figure}[t]
\centering
\begin{tikzpicture}[>=Stealth,scale=0.95]
 % left: compact circle direction with two opposite charges
 \draw[thick] (0,0) circle (1.5);
 \fill[tierC] (-1.5,0) circle (0.11) node[left] {\small $+$};
 \fill[tierL] (1.5,0) circle (0.11) node[right] {\small $-$};
 \foreach \a in {150,120,90,60,30} {\draw[->,thin] (\a+8:1.5) arc (\a+8:\a-8:1.5);}
 \foreach \a in {210,240,270,300,330} {\draw[->,thin] (\a-8:1.5) arc (\a-8:\a+8:1.5);}
 \node at (0,-2.1) {\small (a) total charge $0$: consistent};
 % right: single charge
 \begin{scope}[xshift=6.2cm]
 \draw[thick] (0,0) circle (1.5);
 \fill[tierC] (-1.5,0) circle (0.11) node[left] {\small $+$};
 \foreach \a in {150,120,90,60,30} {\draw[->,thin] (\a+8:1.5) arc (\a+8:\a-8:1.5);}
 \foreach \a in {210,240,270,300,330} {\draw[->,thin] (\a-8:1.5) arc (\a-8:\a+8:1.5);}
 \node[right] at (1.55,0) {\small ?};
 \node at (0,-2.1) {\small (b) total charge $\ne0$: no solution};
 \end{scope}
\end{tikzpicture}
\caption{Gauss's law on a compact space, drawn for a circle. The flux leaving a positive
charge in both directions must be absorbed by a negative charge (a). With a single charge
(b) the two streams of flux meet head-on at the antipodal point and
\eqref{eq:tadpole:gausslaw} has no solution.}
\label{fig:tadpole:gauss}
\end{figure}

\begin{physicsbox}[Physical picture: flux lines have no place to go]
Polchinski's phrase for \eqref{eq:tadpole:zero} is that with a nonzero total
Ramond--Ramond source ``the field equations are inconsistent, because R-R flux lines have no
place to go in the compact space'' \source{Pol Lecture~3, ll.~1987--1989}. He contrasts this with a
nonzero total \emph{tension}, which is only a source for the graviton and dilaton and makes
those fields time dependent: an uncancelled tension changes the solution, an uncancelled
charge means there is no solution \source{Pol Lecture~3, ll.~1985--1989}.
\end{physicsbox}

\subsection{Why the word ``tadpole''}
In perturbation theory, a term in the effective action that is \emph{linear} in a field
$\phi$, $\delta S=\int J\,\phi$, produces a one-point function: a Feynman diagram with a
single external leg ending in a blob, which looks like a tadpole. A linear term means that
$\phi=0$ does not solve the equations of motion. If $\phi$ has a potential, the vacuum
merely shifts. For a gauge potential $C$ whose constant mode
has no potential at all, the equation of motion obtained by varying that constant mode is
simply $\int J=0$, and if the total source $\int J$ does not vanish there is no solution.
This is \eqref{eq:tadpole:zero} in the language of diagrams.\index{tadpole}
Sethi, Vafa and Witten introduce the subject in exactly these terms: a one-loop interaction
``results in a one-point function for the $B$ field of the effective two-dimensional theory,
destabilizing the vacuum'' \source{SVW §1, ll.~45--64}.

\subsection{A first example: orientifold planes and sixteen D-branes}
The oldest stringy instance is the type~I string and its T-duals (\cref{ch:dbranes}).
T-dualizing type~I on $9-p$ circles produces a torus $T^{9-p}$ divided by the reflection of
all its coordinates, combined with worldsheet orientation reversal. The reflection has
$2^{9-p}$ fixed points on the torus; at each sits an \emph{orientifold plane}
(O$p$-plane),\index{orientifold plane} a $p$-dimensional object which is not dynamical but
carries tension and Ramond--Ramond charge. In terms of the charge $\mu_p$ of a D$p$-brane the
result is \source{Pol Lecture~3, eq.~(93), ll.~1968--1984}
\begin{equation}\label{eq:tadpole:oplane}
 \mu_p'=-2^{\,p-5}\mu_p\quad\text{per O$p$-plane},\qquad
 2^{\,9-p}\times\bigl(-2^{\,p-5}\bigr)\mu_p=-16\,\mu_p\quad\text{in total}
\end{equation}
(for the orthogonal projection; the symplectic projection has the opposite sign). The total is
independent of $p$, as it must be, since T-duality cannot change whether the theory is
consistent. Gauss's law \eqref{eq:tadpole:zero} on the compact torus then demands exactly
$16$ D$p$-branes, ``giving the T-dual of $SO(32)$'' \source{Pol Lecture~3, ll.~1989--1990}.
\Cref{tab:tadpole:oplanes} lists the cases used in this chapter.

\begin{table}[t]
\centering
\begin{tabular}{@{}ccccc@{}}
\toprule
$p$ & compact space & number of O$p$-planes & charge of each (units of $\mu_p$) & total\\
\midrule
$9$ & --- & $1$ & $-16$ & $-16$\\
$7$ & $T^2/\Z_2$ & $4$ & $-4$ & $-16$\\
$5$ & $T^4/\Z_2$ & $16$ & $-1$ & $-16$\\
$3$ & $T^6/\Z_2$ & $64$ & $-\tfrac14$ & $-16$\\
\bottomrule
\end{tabular}
\caption{Orientifold-plane charges from \eqref{eq:tadpole:oplane}. In each row sixteen
D$p$-branes cancel the total. The last row agrees with the statement that $T^6/\Z_2$ is a
compact solution ``with 16 D3-branes and 64 O3-planes'', so that the O3 tension is
$-\tfrac14T_3$ \source{GKP §2.1, ll.~419--422}.}
\label{tab:tadpole:oplanes}
\end{table}

\begin{pitfallbox}[Common pitfall: counting branes and their images]
In the covering torus every D-brane away from a fixed point has a mirror image under the
reflection. Some authors count the brane and its image separately (``32 D-branes''), assign
each half the charge, and double the orientifold charge accordingly; others rescale all
charges by a common factor. Only the \emph{ratio} of plane charge to brane charge,
$-2^{p-5}$, and the vanishing of the total are convention independent. The Lean libraries contain
three different bookkeepings of this cancellation (\cref{sec:tadpole:lean}); none is wrong
as arithmetic, but they cannot be compared digit by digit.
\end{pitfallbox}

\section{Curvature as a charge: the origin of \texorpdfstring{$\chi/24$}{chi/24}}
\label{sec:tadpole:chi24}

\subsection{The interaction \texorpdfstring{$C\wedge X_8$}{C X8}}
So far the sources were branes and planes. The new feature of compactifications to two, three
and four dimensions is that the \emph{curvature of the internal space itself} is a source.
In type IIA string theory there is a one-loop interaction
\begin{equation}\label{eq:tadpole:BX8}
 \delta S=-\int B\wedge X_8(R),
\end{equation}
where $B$ is the Neveu--Schwarz two-form, the field that couples to the string, and $X_8(R)$
is an eight-form built as a quartic polynomial in the curvature
\source{SVW §1, eq.~(1.1), ll.~45--54}. The same term appears in M-theory with $B$ replaced
by the three-form $C$ that couples to the membrane, and in F-theory with the four-form that
couples to the D3-brane \source{SVW §1, ll.~64--71}. In eleven-dimensional language it is an
eight-derivative correction to supergravity, suppressed by six powers of the Planck mass,
of the form $-\int C\wedge X_8(R)$ \source{DRS §2.1, eq.~(2.1), ll.~167--180}.

Compactify on an eight-manifold $N$, so that spacetime is $\R^{n}\times N$ with $n=2,3,4$ for
type IIA, M-theory and F-theory respectively. If $C$ is constant along $N$ and has its legs
along $\R^n$, \eqref{eq:tadpole:BX8} becomes $-\bigl(\int_N X_8\bigr)\int_{\R^n}C$: a term
linear in $C$, a tadpole, with coefficient
\begin{equation}\label{eq:tadpole:I}
 I\;=\;-\int_N X_8 \;=\;\frac{\chi(N)}{24}
 \qquad\text{\source{SVW §1, eq.~(1.2) and ll.~93--107}.}
\end{equation}
The first equality is the definition of $I$. The second holds for a Calabi--Yau fourfold;\index{Calabi--Yau fourfold} SVW
obtain it by observing that the one-loop computation reduces to the worldsheet index
$\Tr(-1)^F$ of the sigma model on $N$, which is the Euler number, with a proportionality
constant fixed in the original calculation. Becker and Becker reach the same conclusion from
the differential-form side: for an eight-manifold admitting a nowhere-vanishing spinor the
Pontryagin numbers obey $p_1^2-4p_2+8\chi=0$, the polynomial $X_8$ is proportional to
$P_1^2-4P_2$, and therefore $\int X_8$ is a fixed multiple of $\chi$
\source{BB §2, eqs.~(2.19)--(2.22), ll.~292--339}. \TL

\begin{pitfallbox}[Common pitfall: normalizations of \texorpdfstring{$X_8$}{X8} and of
\texorpdfstring{$G$}{G}]
The three papers normalize differently. BB write $\int_{K^8}X_8=-\chi/\bigl(4!\,(2\pi)^4\bigr)$
and a flux constraint $\int F\wedge F+\chi/12=0$ for their field strength $F$
\source{BB eqs.~(2.22), (2.58), ll.~328--339, 706--713}; SVW use $I=-\int X_8=\chi/24$; DRS
absorb a sign and write the source as $+4\pi^2X_8$ in \eqref{eq:tadpole:warp} below. This
book follows DRS and GVW: the integrally quantized object is $G/2\pi$, and the tadpole
condition is \eqref{eq:tadpole:master}. When you read another paper, fix the dictionary by
the one convention-independent statement: with no flux, the number of branes is $\chi/24$.
\end{pitfallbox}

\subsection{Cancelling the tadpole with branes}
The membrane\index{membrane} couples to $C$ by $\int_{\R^3\times\{y\}}C$ with unit charge. Placing $n$
membranes at points $y_1,\dots,y_n$ of $N$, stretched along $\R^3$, adds $n$ to the
coefficient of the term linear in $C$, with the sign such that branes (not anti-branes)
compensate a positive $I$. So the simplest cure for the tadpole is to ``include in the vacuum
$I$ strings, membranes, and three-branes in Type IIA, M-theory, and F-theory, respectively''
\source{SVW §1, ll.~74--78}. Two obstructions are immediately visible
\source{SVW §1, ll.~79--92}:
\begin{enumerate}[leftmargin=2em]
\item $I$ may be \emph{fractional}. Branes carry integer charge, so nothing can be done with
 branes alone unless $\chi/24\in\Z$.
\item $I$ may be a \emph{negative} integer. Then anti-branes are needed. The vacuum is free
 of tadpoles, but supersymmetry is completely broken: the fourfold preserves supersymmetries
 of one eight-dimensional chirality, a brane preserves those of the same chirality by
 definition, and an anti-brane preserves only the opposite ones.
\end{enumerate}
Hence, with branes alone, a supersymmetric tadpole-free vacuum exists if and only if
$\chi/24$ is a non-negative integer. \TL

\subsection{Cancelling the tadpole with flux}
There is a second contribution. Eleven-dimensional supergravity contains the Chern--Simons
interaction $\int C\wedge G\wedge G$ with $G=\dd C$ \source{DRS §1, eq.~(1.1), ll.~59--68}.
If $G$ has an expectation value along $N$, this term too is linear in the components of $C$
along $\R^3$: a flux on the internal space carries membrane charge. The cleanest way to see
all three contributions at once is the equation for the warp factor. With flux the metric is
no longer a product but
\begin{equation}\label{eq:tadpole:metric}
 \dd s^2=\ee^{-\phi(y)}\,\eta_{\mu\nu}\,\dd x^\mu\dd x^\nu
 +\ee^{\frac12\phi(y)}\,g_{a\bar b}\,\dd y^a\dd y^{\bar b},
\end{equation}
with $g$ a Ricci-flat K\"ahler metric on the fourfold, and the warp factor\index{warp factor} obeys
\source{DRS §2.1, eqs.~(2.2), (2.5), ll.~180--228}
\begin{equation}\label{eq:tadpole:warp}
 \Delta\bigl(\ee^{\frac32\phi}\bigr)=\star\Bigl\{4\pi^2X_8-\tfrac12\,G\wedge G
 -4\pi^2\sum_{i=1}^{n}\delta^8(y-y_i)\Bigr\}.
\end{equation}
Here $\Delta$ and $\star$ are the Laplacian and Hodge star of $g$, and the $\delta^8$ are
eight-forms of unit integral. Equation \eqref{eq:tadpole:warp} is Gauss's law
\eqref{eq:tadpole:gausslaw} for the electric component
$G_{\mu\nu\rho a}=\epsilon_{\mu\nu\rho}\partial_a\ee^{-\frac32\phi}$ of the four-form: the
warp factor plays the role of the electrostatic potential. Now integrate over the compact
fourfold. The left-hand side is the integral of a Laplacian of a globally defined function
and vanishes. Therefore
\begin{equation}\label{eq:tadpole:integrated}
 0=4\pi^2\int_N X_8-\frac12\int_N G\wedge G-4\pi^2 n
 \quad\Longrightarrow\quad
 \int_NX_8=n+\frac{1}{8\pi^2}\int_NG\wedge G .
\end{equation}
Setting $G=0$ and comparing with the brane-only result $n=\chi/24$ of
\eqref{eq:tadpole:I} fixes $\int_NX_8=\chi/24$ in this normalization. Since
$8\pi^2=2(2\pi)^2$, we obtain
\begin{equation}\label{eq:tadpole:drs}
 \frac{\chi}{24}=\frac{1}{8\pi^2}\int_N G\wedge G+n
 =\frac12\int_N\frac{G}{2\pi}\wedge\frac{G}{2\pi}+n ,
\end{equation}
which is \eqref{eq:tadpole:master} \source{DRS §1, eq.~(1.2), ll.~70--80; GVW §2.1,
eq.~(2.1), ll.~94--110}. GVW allow anti-membranes as well and let $n$ be the difference of
the numbers of membranes and anti-membranes, ``a positive or negative integer''. \TL

\subsection{The flux term is non-negative}
Supersymmetry restricts the internal flux to be a $(2,2)$-form that is \emph{primitive},
$J\wedge G=0$ with $J$ the K\"ahler form, and these conditions imply that $G$ is self-dual,\index{self-dual flux}
$G=\star G$ \source{DRS §2.1, eqs.~(2.3), (2.8)--(2.10), ll.~193--263}.\index{primitive form}
\index{G-flux} For any real form, $G\wedge\star G=|G|^2\,\mathrm{vol}$ with $|G|^2\ge0$
pointwise. Hence
\begin{equation}\label{eq:tadpole:pos}
 \int_NG\wedge G=\int_NG\wedge\star G=\int_N|G|^2\,\mathrm{vol}\;\ge\;0,
 \qquad\text{with equality only for }G=0 .
\end{equation}
Combined with $n\ge0$ (no anti-branes, for supersymmetry), \eqref{eq:tadpole:drs} becomes a
\emph{budget}: two non-negative quantities add up to the fixed number $\chi/24$. In
particular
\begin{equation}\label{eq:tadpole:budget}
 0\le n\le\frac{\chi}{24},\qquad 0\le\frac12\int_N\frac{G}{2\pi}\wedge\frac{G}{2\pi}\le
 \frac{\chi}{24},\qquad G=0\iff n=\frac{\chi}{24}.
\end{equation}
DRS state the middle inequality for $K3\times K3$ as
$\frac12\int\frac G{2\pi}\wedge\frac G{2\pi}\le24$ \source{DRS §3.1, ll.~640--649}. The
elementary step from \eqref{eq:tadpole:drs} to \eqref{eq:tadpole:budget}, for natural numbers
adding up to $24$, is the Lean theorem \lean{tadpole\_budget}
(\cref{sec:tadpole:lean}).\index{tadpole!budget}

\begin{physicsbox}[Physical picture: an electrostatics problem with three kinds of charge]
Read \eqref{eq:tadpole:warp} as electrostatics on the fourfold. The curvature is a smooth
background charge density of total charge $+\chi/24$ (in units where a membrane has charge
$1$ and we have moved everything to one side). Membranes are point charges; self-dual flux
is a smooth charge density of the same sign as the membranes. The warp factor is the
potential they generate together, and it exists only if the total vanishes. The only
situation in which one may forget all this is the non-compact one: on $\R^8$, or on a local
model of a singularity, ``flux can escape to infinity, and hence [the tadpole condition] no
longer holds''; the flux measured at infinity appears as an extra term
\source{GVW §2.2, ll.~380--399}.
\end{physicsbox}

\section{How large is \texorpdfstring{$\chi/24$}{chi/24}?}\label{sec:tadpole:euler}

The left-hand side of \eqref{eq:tadpole:master} is pure topology. We compute it in three
ways for $K3\times K3$, and then ask what can be said in general.

\subsection{Three routes to \texorpdfstring{$\chi(K3\times K3)=576$}{chi(K3 x K3) = 576}}
\emph{Route 1: K\"unneth.}\index{K\"unneth formula} The cohomology of a product is the tensor
product of the cohomologies, so Poincar\'e polynomials multiply. With
$P_{K3}(t)=1+22t^2+t^4$,
\begin{equation}\label{eq:tadpole:poincare}
 P_{K3\times K3}(t)=\bigl(1+22t^2+t^4\bigr)^2=1+44t^2+486t^4+44t^6+t^8 ,
\end{equation}
so the Betti numbers are $(1,0,44,0,486,0,44,0,1)$ and
$\chi=P(-1)=1+44+486+44+1=576=24^2$. The Euler number is multiplicative because
$P_{X\times Y}(-1)=P_X(-1)P_Y(-1)$. The middle Betti number decomposes as
\begin{equation}\label{eq:tadpole:b4}
 b_4=486=\underbrace{1\cdot1}_{H^0\otimes H^4}+\underbrace{22\cdot22}_{H^2\otimes H^2}
 +\underbrace{1\cdot1}_{H^4\otimes H^0},
\end{equation}
a decomposition we shall need in \cref{sec:tadpole:integrality}: $G$-flux lives in $H^4$.

\emph{Route 2: Hodge numbers.} The Hodge numbers of a product are given by the same rule,
$h^{p,q}(X\times Y)=\sum_{r,s}h^{r,s}(X)\,h^{p-r,q-s}(Y)$. From the K3 diamond one finds
the $5\times5$ table
\begin{equation}\label{eq:tadpole:hodge}
 \bigl(h^{p,q}(K3\times K3)\bigr)=
 \begin{pmatrix}1&0&2&0&1\\0&40&0&40&0\\2&0&404&0&2\\0&40&0&40&0\\1&0&2&0&1\end{pmatrix},
 \qquad \chi=\sum_{p,q}(-1)^{p+q}h^{p,q}=576 .
\end{equation}
For instance $h^{2,2}=1\cdot1+20\cdot20+1\cdot1+2\cdot(1\cdot1)=404$: the last term counts
$h^{2,0}h^{0,2}+h^{0,2}h^{2,0}$. Every nonzero entry has $p+q$ even, so all signs are $+$ and
the entries add up to $576$. Note $h^{2,0}=2$: the product has two holomorphic two-forms,
one from each factor. Its holonomy is $SU(2)\times SU(2)$, smaller than $SU(4)$, and M-theory
on it has $N=4$ rather than $N=2$ supersymmetry in three dimensions
\source{DRS §3.1, ll.~638--640}.

\emph{Route 3: Chern classes.} For a complex fourfold $\chi=\int c_4$. The tangent bundle of
a product is the sum of the pull-backs, so total Chern classes multiply:
$c(K3\times K3)=(1+c_2^{(1)})(1+c_2^{(2)})$, where $c_2^{(i)}$ is the second Chern class of
the $i$-th factor, with $\int_{K3}c_2=24$ and $c_1=0$. Hence
\begin{equation}\label{eq:tadpole:chern}
 c_2=c_2^{(1)}+c_2^{(2)},\qquad c_4=c_2^{(1)}c_2^{(2)},\qquad
 \int c_4=24\cdot24=576,\qquad \int c_2^2=2\cdot24\cdot24=1152 .
\end{equation}
All three routes give $\chi/24=24$, the value quoted by SVW and DRS
\source{SVW §2.2, ll.~252--255; §2.4, l.~535; DRS §3.1, l.~640}. The arithmetic
$576=24\cdot24$ and $576/24=24$ with remainder $0$ is the Lean theorem
\lean{k3k3\_anomaly}. \TA{} for the arithmetic, \TL{} for the K\"unneth formula.

\subsection{The general fourfold: \texorpdfstring{$\chi$}{chi} is divisible by six}
Is $\chi/24$ always an integer? No, and the argument of SVW shows exactly how far one can
go. Let $N$ be a Calabi--Yau fourfold and
$\chi_p=\sum_{q=0}^{4}(-1)^qh^{p,q}$ the index of the Dolbeault operator on $(p,0)$-forms.
The index theorem expresses $\chi_p$ through Chern numbers; since $c_1=0$ only
$\int c_2^2$ and $\int c_4=\chi$ can appear \source{SVW §2.1, eq.~(2.1), ll.~125--155}:
\begin{equation}\label{eq:tadpole:todd}
 \chi_0=\frac{1}{720}\int\bigl(3c_2^2-c_4\bigr),\qquad
 \chi_1=\frac{1}{180}\int\bigl(3c_2^2-31c_4\bigr),\qquad
 \chi_2=\frac{1}{120}\int\bigl(3c_2^2+79c_4\bigr).
\end{equation}
Eliminate $\int c_2^2$ with the first equation, $\int c_2^2=240\chi_0+\chi/3$:
\begin{equation}\label{eq:tadpole:div6}
 \chi_1=\frac{720\chi_0+\chi-31\chi}{180}=4\chi_0-\frac{\chi}{6},\qquad
 \chi_2=\frac{720\chi_0+\chi+79\chi}{120}=6\chi_0+\frac{2\chi}{3}.
\end{equation}
For holonomy exactly $SU(4)$ one has $h^{0,q}=(1,0,0,0,1)$, so $\chi_0=2$, and
\eqref{eq:tadpole:div6} becomes the formulas of SVW, $\int c_2^2=480+\chi/3$,
$\chi_1=8-\chi/6$ and $\chi_2=12+2\chi/3$ \source{SVW §2.1, ll.~158--204}. Since $\chi_1$ and
$\chi_0$ are integers, \emph{$\chi$ is divisible by $6$}; the possible fractional parts of
$\chi/24$ are $0,\tfrac14,\tfrac12,\tfrac34$ and ``the obstruction takes values in $\Z_4$''
\source{SVW §2.1, ll.~121--124}.\index{Euler characteristic!divisibility by six} \TL

\begin{example}[Checking \eqref{eq:tadpole:div6} on $K3\times K3$]\label{ex:tadpole:check}
Here $\chi_0=1-0+2-0+1=4$ from the first row of \eqref{eq:tadpole:hodge}, not $2$, because the
holonomy is smaller than $SU(4)$. The second row gives $\chi_1=-40-40=-80$ and the third
$\chi_2=2+404+2=408$. The formulas predict $4\cdot4-576/6=16-96=-80$ and
$6\cdot4+2\cdot576/3=24+384=408$. The Chern numbers \eqref{eq:tadpole:chern} give
$\chi_0=(3\cdot1152-576)/720=2880/720=4$. All agree (symbolic check with sympy).
\end{example}

\begin{example}[A fractional tadpole: the sextic fourfold]\label{ex:tadpole:sextic}
Let $N\subset\mathbb P^5$ be a smooth hypersurface of degree six, a Calabi--Yau fourfold.
By adjunction its total Chern class is $c(N)=(1+H)^6/(1+6H)$ with $H$ the hyperplane class.
Expanding,
\[
 \frac{(1+H)^6}{1+6H}=1+0\cdot H+15H^2-70H^3+435H^4+\dots
\]
The vanishing coefficient of $H$ confirms $c_1=0$. Since $\int_NH^4=6$ (a degree-six
hypersurface meets a line in six points), $\chi=6\cdot435=2610$. This is divisible by $6$
($2610=6\cdot435$) but $2610=24\cdot108+18$, so
\[
 \frac{\chi}{24}=108\tfrac34 ,
\]
as stated in \source{SVW §2.1, ll.~205--213}. No number of membranes cancels this tadpole.
DRS point out the way out: ``for cases where $\chi/24$ is not integral, $G$-flux is actually
required'' \source{DRS §1, ll.~70--72}. Such a flux cannot be an integral class $\xi$: then
$\tfrac12\int\xi\wedge\xi\in\tfrac12\Z$ could never supply the fractional part $\tfrac34$.
It lies in the shifted lattice described in \cref{sec:tadpole:integrality}.
\end{example}

\subsection{Elliptic fibrations, and a negative example}
For F-theory one needs $N$ elliptically fibred over a threefold $B$, with a section and a
smooth Weierstrass model. Then $\chi$ can be written in terms of the Chern classes $c_i$ of
the base \source{SVW §2.2, eqs.~(2.3), (2.10)--(2.12), ll.~258--386}:
\begin{equation}\label{eq:tadpole:elliptic}
 \chi(N)=12\int_Bc_1\bigl(c_2+30c_1^2\bigr),\qquad
 \frac{\chi}{24}=\int_B\Bigl(15c_1^3+\tfrac12c_1c_2\Bigr).
\end{equation}
Two checks. For $B=K3\times\mathbb P^1$ the fourfold is $K3\times K3$ with one factor
elliptic; $c_1(B)$ is twice the point class of $\mathbb P^1$, so $c_1^3=0$ and
$\int c_1c_2=2\cdot24=48$, giving $\chi/24=24$. For $B=\mathbb P^3$, $c_1=4H$ and $c_2=6H^2$,
so $\int c_1c_2=24$, $\int c_1^3=64$ and $\chi/24=15\cdot64+12=972$, the value quoted by SVW.
When $N$ has full $SU(4)$ holonomy the arithmetic genus of $B$ is $\int_Bc_1c_2/24=1$, and
\eqref{eq:tadpole:elliptic} becomes $\chi/24=12+15\int_Bc_1^3$, a positive integer: in this
class supersymmetric tadpole cancellation by branes is always possible
\source{SVW §2.2, eq.~(2.4), ll.~270--291}. \TL

The sign obstruction is rarer. SVW find it in a conformal field theory rather than a smooth
manifold: the orbifold $T^8/(\Z_2\times\Z_2)$ is a singular point of $K3\times K3$, but
\emph{with discrete torsion}\index{discrete torsion} its Hodge diamond becomes
\begin{equation}\label{eq:tadpole:torsion}
 \begin{pmatrix}1&0&2&0&1\\0&8&64&8&0\\2&64&20&64&2\\0&8&64&8&0\\1&0&2&0&1\end{pmatrix},
 \qquad \chi=4-48-104-48+4=-192=-8\cdot24 ,
\end{equation}
where we summed the rows with signs $(-1)^{p+q}$. Thus $\chi/24=-8$: eight anti-branes are
needed and supersymmetry is broken \source{SVW §2.3, ll.~429--515}. The large odd entries
$h^{2,1}=h^{1,2}=64$ are what drives $\chi$ negative.

\section{Flux quantization and why \texorpdfstring{$\tfrac12\int G\wedge G$}{half of G wedge G}
is an integer}\label{sec:tadpole:integrality}

\subsection{Dirac quantization of \texorpdfstring{$G$}{G}}
The flux term in \eqref{eq:tadpole:master} must be an integer whenever $\chi/24$ is, for
otherwise no integer $n$ solves the equation. Where does this come from? The membrane
couples to $C$ as a charged particle couples to a gauge potential, and the usual Dirac
argument (the phase $\exp(\ii\int C)$ must be well defined when the membrane worldvolume is
deformed across a four-cycle) requires the periods of $G/2\pi$ over four-cycles to be
integers:\index{flux quantization}
\begin{equation}\label{eq:tadpole:dirac}
 \xi:=\Bigl[\frac{G}{2\pi}\Bigr]\in H^4(N;\Z)\qquad\text{(first approximation).}
\end{equation}
The precise statement carries a shift. The periods of $G/2\pi$ coincide with those of
$c_2(N)/2$ modulo integers \source{GVW §5, ll.~1993--1996}, so that $\xi$ lives in a
translate of the lattice $H^4(N;\Z)$; this shift ``arises precisely when the intersection
form on $H^4(Y;\Z)$ is not even'', and with the shifted condition the right-hand side of the
tadpole formula is always integral \source{GVW §2.1, ll.~112--125}. DRS phrase the same fact
as: $[G/\pi]$ is always integral, and ``if $\chi/24\in\Z$ then $[G/2\pi]$ is an integer
cohomology class'' \source{DRS §2.1, ll.~230--234}. \TL

For $K3\times K3$ there is no shift, and we can see this directly. By
\eqref{eq:tadpole:chern}, $c_2/2=\tfrac12c_2^{(1)}+\tfrac12c_2^{(2)}$. The class
$c_2^{(1)}$ is $24$ times the generator of $H^4(K3_1;\Z)$, so $c_2/2$ is $12$ times a
generator plus $12$ times another: an integral class. Its periods are integers, and
\eqref{eq:tadpole:dirac} holds as written. What remains is to understand why \emph{half} the
self-intersection of an integral class is an integer. That is a property of the lattice.

\subsection{The lattice \texorpdfstring{$H^4(K3\times K3;\Z)$}{H4(K3 x K3; Z)}}
Let $\Lambda=H^2(K3;\Z)\cong3U\oplus2E_8(-1)$ with Gram matrix $L$ in a basis $e_1,\dots,
e_{22}$, so $L_{ik}=\int_{K3}e_i\wedge e_k$. The K\"unneth decomposition
\eqref{eq:tadpole:b4} holds over $\Z$ because the cohomology of K3 has no torsion:
\begin{equation}\label{eq:tadpole:H4}
 H^4(K3_1\times K3_2;\Z)\;=\;\underbrace{\Z\,[K3_1\times\mathrm{pt}]^\vee\oplus
 \Z\,[\mathrm{pt}\times K3_2]^\vee}_{\text{rank }2}\;\oplus\;
 \underbrace{\Lambda_1\otimes\Lambda_2}_{\text{rank }484}.
\end{equation}
The two classes in the first summand are the volume forms of the two factors; each squares
to zero (a volume form of a four-manifold wedged with itself vanishes) and their product
integrates to $1$, so they span a hyperbolic plane $U$. On the second summand take the basis
$e_i\otimes f_j$, represented by $\pi_1^*e_i\wedge\pi_2^*f_j$. Two-forms commute, so
\begin{equation}\label{eq:tadpole:kron}
 \int_{K3_1\times K3_2}(e_i\wedge f_j)\wedge(e_k\wedge f_l)
 =\int_{K3_1}e_i\wedge e_k\int_{K3_2}f_j\wedge f_l=L_{ik}\,L_{jl} .
\end{equation}
The Gram matrix of $\Lambda_1\otimes\Lambda_2$ is therefore the \emph{Kronecker product}
$L\otimes L$, a $484\times484$ integer matrix with entries
$(L\otimes L)_{(i,j),(k,l)}=L_{ik}L_{jl}$.\index{Kronecker product} A flux in this summand is
$\xi=\sum_{i,j}x_{ij}\,e_i\otimes f_j$ with integers $x_{ij}$, and
\begin{equation}\label{eq:tadpole:Q}
 \int\xi\wedge\xi=Q(x):=\sum_{i,j,k,l}x_{ij}\,L_{ik}L_{jl}\,x_{kl}.
\end{equation}
For a decomposable flux $\xi=a\otimes b$, i.e.\ $x_{ij}=a_ib_j$, this factorizes:
$Q=(a\cdot a)(b\cdot b)$, the product of the two self-intersections on the two K3 surfaces.

\begin{proposition}[Evenness of the flux lattice]\label{prop:tadpole:even}
Let $L_1$, $L_2$ be symmetric integer matrices and suppose every diagonal entry of $L_1$ is
even. Then $x^{\mathsf T}(L_1\otimes L_2)\,x$ is even for every integer vector $x$. Consequently
$\tfrac12\int\xi\wedge\xi\in\Z$ for every $\xi\in\Lambda_1\otimes\Lambda_2$.
\end{proposition}
\begin{proof}
Write $M=L_1\otimes L_2$ and label its rows by pairs $p=(i,j)$. $M$ is symmetric because
$M_{(i,j),(k,l)}=(L_1)_{ik}(L_2)_{jl}=(L_1)_{ki}(L_2)_{lj}=M_{(k,l),(i,j)}$. Its diagonal is
$M_{pp}=(L_1)_{ii}(L_2)_{jj}$, a product with an even factor, hence even. For any symmetric
integer matrix,
\[
 x^{\mathsf T}Mx=\sum_pM_{pp}\,x_p^2+2\sum_{p<q}M_{pq}\,x_px_q ,
\]
after choosing any total order on the labels. The second sum is visibly even and the first
is even because each $M_{pp}$ is.
\end{proof}
Two remarks. First, \emph{nothing is assumed about $L_2$} beyond symmetry: one even factor
suffices, and unimodularity plays no role. Second, the statement is about all of
$\Lambda_1\otimes\Lambda_2$, not only about decomposable elements, for which it would be
trivial. The proposition is the Lean theorem \lean{kronForm\_even}, and its corollary is
\lean{flux\_half\_selfIntersection\_integral}. \TA

Adding the hyperbolic plane of \eqref{eq:tadpole:H4}, which is also even, the whole lattice
$H^4(K3\times K3;\Z)\cong U\oplus(\Lambda\otimes\Lambda)$ is even, consistently with the
absence of a shift in the quantization condition. Its signature follows from that of
$\Lambda$, $(3,19)$: a tensor product of a positive with a positive or a negative with a
negative direction is positive, so
\begin{equation}\label{eq:tadpole:sig}
\begin{aligned}
 (n_+,n_-)&=\bigl(3\cdot3+19\cdot19+1,\;2\cdot3\cdot19+1\bigr)=(371,115),\\
 371+115&=486,\qquad 371-115=256=(-16)^2 ,
\end{aligned}
\end{equation}
the last equality being the multiplicativity of the signature, $\sigma(K3)=-16$. A
supersymmetric flux is self-dual and so lies in the $371$-dimensional positive part, in
agreement with \eqref{eq:tadpole:pos}.

\begin{pitfallbox}[Common pitfall: ``the form is indefinite, so the flux term can be negative'']
The lattice $H^4$ is indefinite, and an arbitrary integral class $\xi$ can have
$\int\xi\wedge\xi<0$. Such a class is not a solution of the equations of motion with
unbroken supersymmetry: it fails self-duality. The positivity \eqref{eq:tadpole:pos} is a
statement about the solutions, the evenness of \cref{prop:tadpole:even} is a statement about
the whole lattice. The tadpole budget uses both.
\end{pitfallbox}

\subsection{Worked example: three ways to spend a budget of 24}\label{sec:tadpole:worked}
We follow \source{DRS §3.1, eqs.~(3.1)--(3.5), ll.~660--728}. Take for each K3 a smooth
hypersurface of degree $(2,2,2)$ in $\mathbb P^1\times\mathbb P^1\times\mathbb P^1$. It has
three natural divisor classes $C_1,C_2,C_3$, obtained by fixing a point in one of the three
$\mathbb P^1$ factors. Two distinct $C_i$ meet where two coordinates are fixed, which leaves
a quadratic equation on the third $\mathbb P^1$: two points. A $C_i$ can be moved off itself,
so $C_i\cdot C_i=0$. The Gram matrix of the sublattice they span is
\begin{equation}\label{eq:tadpole:drsgram}
 L_{\mathrm{DRS}}=\begin{pmatrix}0&2&2\\2&0&2\\2&2&0\end{pmatrix},\qquad\det L_{\mathrm{DRS}}=16 .
\end{equation}
It is symmetric with even (zero) diagonal, so \cref{prop:tadpole:even} applies to it although
it is far from unimodular. Take $J=C_1+C_2+C_3$ as K\"ahler class, and
\[
 \alpha=C_1-C_2,\qquad\beta=C_1-C_3 .
\]
With $L_{\mathrm{DRS}}$ one computes
$\alpha\cdot\alpha=0+0-2\cdot2=-4$, likewise $\beta\cdot\beta=-4$,
$\alpha\cdot\beta=C_1^2-C_1C_3-C_2C_1+C_2C_3=0-2-2+2=-2$, and
$J\cdot\alpha=(2+2)-(2+2)=0=J\cdot\beta$. So $\alpha,\beta$ are primitive $(1,1)$ classes.
On a K\"ahler surface a primitive $(1,1)$-form is anti-self-dual (the two-form computation in
\source{DRS §2.1, eq.~(2.7), ll.~240--241}), and on a product of two four-manifolds $\star(a\wedge b)=\star_1a\wedge\star_2b$ for
two-forms $a,b$ pulled back from the factors, so the wedge product of two anti-self-dual
two-forms is self-dual: $(-a)\wedge(-b)=a\wedge b$. With subscripts for the
two K3 factors, three flux choices are:
\begin{center}
\begin{tabular}{@{}lccc@{}}
\toprule
flux $\xi=G/2\pi$ & $\int\xi\wedge\xi$ & $\tfrac12\int\xi\wedge\xi$ & membranes $n$\\
\midrule
$0$ & $0$ & $0$ & $24$\\
$\alpha_1\wedge\alpha_2$ & $(-4)(-4)=16$ & $8$ & $16$\\
$(\alpha_1+\beta_1)\wedge\alpha_2$ & $(-12)(-4)=48$ & $24$ & $0$\\
\bottomrule
\end{tabular}
\end{center}
using $(\alpha+\beta)^2=-4-4+2(-2)=-12$. The second and third rows are DRS's equations (3.4)
and (3.5): ``we can place 16 membranes'' or ``cancel the anomaly completely without branes''.
A non-decomposable flux is just as easy with \eqref{eq:tadpole:Q}: for
$\xi=\alpha_1\wedge\alpha_2+\beta_1\wedge\beta_2$,
\[
 \int\xi\wedge\xi=(\alpha\cdot\alpha)^2+(\beta\cdot\beta)^2+2(\alpha\cdot\beta)^2=16+16+8=40,
 \qquad n=24-20=4 .
\]
All self-intersections are even, as \cref{prop:tadpole:even} guarantees, and all four rows
satisfy \eqref{eq:tadpole:budget}. \Cref{fig:tadpole:budget} displays the result.

\begin{figure}[t]
\centering
\begin{tikzpicture}[x=0.42cm,y=0.62cm]
 \foreach \y/\f/\lab in {3/0/{$\xi=0$},2/8/{$\alpha_1\wedge\alpha_2$},
   1/20/{$\alpha_1\wedge\alpha_2+\beta_1\wedge\beta_2$},0/24/{$(\alpha_1+\beta_1)\wedge\alpha_2$}} {
   \fill[tierL!45] (0,\y) rectangle (\f,\y+0.6);
   \fill[tierC!35] (\f,\y) rectangle (24,\y+0.6);
   \draw (0,\y) rectangle (24,\y+0.6);
   \node[left] at (0,\y+0.3) {\small \lab};
 }
 \node at (4,3.3) {\small $n=24$};
 \node at (4,2.3) {\small $8$}; \node at (16,2.3) {\small $n=16$};
 \node at (10,1.3) {\small $20$}; \node at (22,1.3) {\small $4$};
 \node at (12,0.3) {\small $24$};
 \draw[->] (0,-0.5) -- (25.5,-0.5) node[right] {\small $\chi/24$};
 \foreach \x in {0,8,16,24} {\draw (\x,-0.4) -- (\x,-0.6) node[below] {\small \x};}
 \fill[tierL!45] (3,-2.3) rectangle (4,-1.8); \node[right] at (4.1,-2.05) {\small flux
   $\tfrac12\int\xi\wedge\xi$};
 \fill[tierC!35] (14,-2.3) rectangle (15,-1.8); \node[right] at (15.1,-2.05) {\small membranes $n$};
\end{tikzpicture}
\caption{The tadpole budget of $K3\times K3$ for the four fluxes of
\cref{sec:tadpole:worked}. Flux and branes are both non-negative and add up to $\chi/24=24$.}
\label{fig:tadpole:budget}
\end{figure}

\begin{physicsbox}[Physical picture: flux freezes moduli]
The second and third rows are possible only on special K3 surfaces. A generic K3 has no
integral $(1,1)$ classes at all (its Picard lattice is trivial, \cref{ch:k3lattice}), and DRS
begin their construction with exactly this remark \source{DRS §3.1, ll.~673--676}. Once an
integral class $\xi$ is chosen, the requirement that it be of type $(2,2)$ and primitive
is a condition on the complex structure and on the K\"ahler class: for non-zero $G$ ``we
must adjust the complex structure of $X$ so that $G^{0,4}=G^{1,3}=0$'', and the K\"ahler
class so that $G$ is primitive \source{GVW §2.4, eqs.~(2.36)--(2.37), ll.~898--911}.
Turning this around, a chosen flux restricts the moduli to
the locus where it is supersymmetric. This is the mechanism of moduli stabilization, the
subject of \cref{ch:stabilization}. Because both terms of the budget are non-negative, GVW note that there are only finitely many
possible choices of $n$ and $G$ at a fixed value of the total \source{GVW §2.2, ll.~401--405}.
\end{physicsbox}

\section{The type IIB picture: \texorpdfstring{$\KT/\Z_2$}{K3 x T2 / Z2} with O7-planes and
D7-branes}\label{sec:tadpole:iib}

\subsection{From M-theory on \texorpdfstring{$K3\times K3$}{K3 x K3} to the orientifold}
F-theory is reached from M-theory by taking one K3 to be elliptically fibred and shrinking
the fibre; M-theory on the fourfold goes over to type IIB on the base $B$, with the complex
structure $\tau$ of the fibre as the varying type IIB coupling
\source{DRS §1, ll.~45--52}.\index{F-theory} Take the elliptic K3 at its orbifold point
$T^4/\Z_2=(T^2_{\text{base}}\times T^2_{\text{fibre}})/\Z_2$. Then the base of the fibration
is $T^2/\Z_2$ and the fourfold $K3\times T^4/\Z_2$ describes type IIB on
\begin{equation}\label{eq:tadpole:orientifold}
 \frac{\KT}{\Omega(-1)^{F_L}\,\Z_2},\qquad \Z_2:\ z^1\mapsto-z^1\ \text{on }T^2 ,
\end{equation}
where $\Omega$ is worldsheet parity and $(-1)^{F_L}$ the left-moving fermion number
\source{DRS §4.1, ll.~1289--1293; §4.2, ll.~1350--1352}. This orientifold ``is a special
point in the moduli space of F theory on $K3\times K3$''. The dictionary between the two
pictures is explicit in DRS's account of the sixteen fixed points
\source{DRS §2.4, ll.~511--541}:\index{orientifold!of K3 x T2}
\begin{itemize}[leftmargin=2em]
\item M-theory on $T^4/\Z_2$ has $22$ gauge fields from reducing $C$ on $H^2(K3;\Z)$: sixteen
 from the twisted sectors (the $A_1$ singularities of \cref{ch:kummer}) and six untwisted.
\item In the F-theory limit the two untwisted forms dual to the fibre and its Hodge dual
 drop out; the element $\Omega(-1)^{F_L}$ acts on the doublet of two-forms $(B,B')$ as
 $-\mathbf 1\in\SLZ$, so only components with one leg on $T^2$ survive.
\item ``The 16 fixed points coalesce into 4 groups of coincident $A_1$ singularities'', one
 group at each of the four fixed points of $z^1\mapsto-z^1$ on $T^2$. The gauge group is
 enhanced $SU(2)^4\to SO(8)$ at each, which in the orientifold arises ``from placing 4
 D7-branes at the location of each O7-plane''. The total gauge group is $SO(8)^4$
 \source{DRS §4.1, ll.~1311--1315}.
\end{itemize}

\subsection{Two tadpoles, not one}
The orientifold \eqref{eq:tadpole:orientifold} has \emph{two} Gauss laws, and it is important
to keep them apart.

\emph{The seven-brane tadpole.}\index{O7-plane}\index{D7-brane} O7-planes and D7-branes fill spacetime and wrap K3; they are
points on $T^2/\Z_2$. They are charged under the Ramond--Ramond eight-form $C_8$, whose
field strength has flux lines spreading in the two transverse directions, the compact
pillow $T^2/\Z_2$. Gauss's law \eqref{eq:tadpole:zero} on the pillow, with the $p=7$ row
of \cref{tab:tadpole:oplanes}, reads
\begin{equation}\label{eq:tadpole:d7}
 \underbrace{4\times(-4)}_{\text{O7-planes}}+\underbrace{16\times(+1)}_{\text{D7-branes}}=0
 \qquad\text{(units of }\mu_7\text{)} .
\end{equation}
With four D7-branes on top of each O7-plane the cancellation is \emph{local}:
$-4+4=0$ at each fixed point, no $C_8$ flux lines run across the pillow, the coupling
$\tau$ is constant, and perturbation theory applies everywhere. This is why the $SO(8)^4$
point is the orientifold limit of F-theory (\cref{fig:tadpole:pillow}). No choice of
three-form flux enters \eqref{eq:tadpole:d7}.

\emph{The three-brane tadpole.}\index{D3-brane}\index{tadpole!three-brane} D3-branes fill spacetime and are points in the full
six-dimensional internal space. They are charged under the four-form $D^+$ of type IIB,
the field into which the membrane's $C_{\mu\nu\rho}$ lifts
\source{DRS §2.4, eq.~(2.27), ll.~479--486}. A $G$-flux with one leg along the elliptic fibre,
$G/2\pi=\dd z\wedge\omega-\dd\bar z\wedge\star\omega$ with $\omega\in H^{1,2}(B)$, lifts to
the Neveu--Schwarz and Ramond--Ramond three-form field strengths
\begin{equation}\label{eq:tadpole:HH}
 H=\omega-\star\omega,\qquad H'=\omega\tau-\star\omega\,\bar\tau ,
\end{equation}
and the tadpole condition \eqref{eq:tadpole:master} becomes
\begin{equation}\label{eq:tadpole:d3}
 \frac{\chi}{24}=n-\int_BH\wedge H' ,\qquad n=\text{number of D3-branes},
\end{equation}
in the normalization of \source{DRS §2.4, eqs.~(2.30)--(2.32), ll.~549--578}; the
combination $-\int H\wedge H'$ is the image of the non-negative M-theory flux term. It
comes from the type IIB Chern--Simons interaction $\int D^+\wedge H\wedge H'$
\source{DRS §1, eq.~(1.3), ll.~92--100}, the descendant of $\int C\wedge G\wedge G$. In the
conventions of GKP the same statement is the integrated Bianchi identity for the five-form,
$(2\kappa_{10}^2T_3)^{-1}\int H_{(3)}\wedge F_{(3)}+Q_3^{\rm loc}=0$, with $Q_3^{\rm loc}$
the D3 charge of all localized sources \source{GKP §2.1, eqs.~(2.23)--(2.24), ll.~483--502}.
For our orientifold the number on the
left of \eqref{eq:tadpole:d3} is $\chi(K3\times K3)/24=24$,
\begin{equation}\label{eq:tadpole:d3k3}
 -\int H\wedge H'+n=24 ,
\end{equation}
which DRS write in its M-theory form $\frac12\int\frac G{2\pi}\wedge\frac G{2\pi}+n=24$ on
$K3\times T^4/\Z_2$ before lifting the flux \source{DRS §4.2, eqs.~(4.4)--(4.5),
ll.~1358--1390}.

\begin{figure}[t]
\centering
\begin{tikzpicture}[scale=1.0]
 \draw[thick,rounded corners=10pt,fill=blue!4] (0,0) rectangle (4.2,2.6);
 \foreach \x/\y in {0.45/0.45,3.75/0.45,0.45/2.15,3.75/2.15} {
   \draw[tierC,very thick] (\x-0.11,\y-0.11) -- (\x+0.11,\y+0.11) (\x-0.11,\y+0.11) -- (\x+0.11,\y-0.11);
   \foreach \k in {0,1,2,3} {\fill[tierL] (\x+0.17+0.11*\k,\y+0.24) circle (0.04);}
 }
 \node at (2.1,1.3) {\small $T^2/\Z_2$};
 \node[right,align=left] at (4.6,2.0) {\small $\times$: O7-plane, $C_8$ charge $-4$};
 \node[right,align=left] at (4.6,1.45) {\small $\bullet$: D7-brane, $C_8$ charge $+1$ each};
 \node[right,align=left] at (4.6,0.9) {\small per corner: $-4+4=0$\quad($SO(8)$)};
 \node[right,align=left] at (4.6,0.35) {\small D3 budget of the whole space: $24$};
\end{tikzpicture}
\caption{Charge bookkeeping on the pillow $T^2/\Z_2$, in units of the D7 charge $\mu_7$.
Each seven-brane also wraps the K3 (not drawn). The seven-brane tadpole cancels corner by
corner; the three-brane tadpole $24$ must still be cancelled by D3-branes and three-form
flux. The geometry of the pillow is described in \cref{ch:k3t2}.}
\label{fig:tadpole:pillow}
\end{figure}

\begin{remark}[Where the $24$ sits in the type IIB picture]\label{rem:tadpole:where}
In M-theory the number $24$ is smeared over the curvature of $K3\times K3$. In type IIB the
internal space $\KT$ has $\chi=0$ (\cref{ch:k3t2}) and cannot be the source. The source is
the curvature of the \emph{seven-brane worldvolumes}: through a Chern--Simons coupling of the
form $\int C_{(4)}\wedge p_1(R)$, ``a D7-brane wrapped on K3 has $-1$ unit of D3 charge'', and
in F-theory the total background charge of all seven-branes is
$Q_3^{\rm D7}=-\chi(X)/24$ \source{GKP §2.1, eqs.~(2.20), (2.22), ll.~428--474}; SVW make the
same observation, the induced charge of a worldvolume $K$ being $c_2(K)/24$
\source{SVW §2.2, ll.~387--399}. Note $\chi(K3)/24=1$. Sixteen D7-branes account for $-16$
of the $-24$. The remaining $-8$ must be carried by the four O7-planes, $-2$ each; this last
division is standard but has not been checked against a source in this book's library.
\end{remark}

\subsection{A chain of dualities as a consistency check}
With no flux, \eqref{eq:tadpole:d3k3} requires $24$ D3-branes. Two T-dualities along the
$T^2$ invert both radii, $R_i\mapsto\ap/R_i$ (\cref{ch:tduality}), map
$\Omega(-1)^{F_L}\Z_2$ to $\Omega$, the O7/D7 system to the O9/D9 system of type~I, and each
D3-brane to a D5-brane wrapped on the dual torus: ``the resulting theory is type I on
$\KT$ with 24 D5-branes wrapping the $T^2$'' \source{DRS §4, ll.~1270--1276; §4.1,
ll.~1298--1311}. An S-duality turns this into the $SO(32)$ heterotic string on $\KT$. SVW use exactly this chain, in the opposite direction, as ``a check on the proportionality
constant in the formula $I=\chi/24$'': F-theory on $K3\times K3$ is dual to the heterotic
string on $\KT$, where, if no gauge bundle is excited on the K3, the worldsheet anomaly is
cancelled by $24$ five-branes wrapped on $T^2$; the eight-dimensional duality turns them into
$24$ three-branes \source{SVW §2.4, ll.~526--535}. \TL


\begin{physicsbox}[Physical picture: what flux buys, and what it costs]
Replacing some of the $24$ D3-branes by three-form flux changes the four-dimensional physics.
DRS's flux \eqref{eq:tadpole:HH} on $\KT/\Z_2$ is built from a class $\alpha\in H^{1,1}(K3)$
and a class $\beta\in H^{2,0}(K3)$; ``if $\beta=0$, the model has $N=2$ supersymmetry'' and
otherwise $N=1$ \source{DRS §4.2, eq.~(4.4), ll.~1374--1383}. The metric is no longer a
product but is warped, conformal to $\KT$ \source{DRS §4, ll.~1268--1270}, and after the two
T-dualities the type~I and heterotic duals live on a space that is complex but not K\"ahler,
with non-zero three-form flux (torsion) \source{DRS §4, ll.~1276--1286}. T-duality is
here the tool that turns a flux compactification one understands into one that would be hard
to find directly.
\end{physicsbox}

\section{What the machine has checked}\label{sec:tadpole:lean}

The physics of this chapter (the interaction $C\wedge X_8$, the identification
$\int X_8=\chi/24$, supersymmetry conditions, dualities) is Tier~L throughout. What Lean
checks is the integer arithmetic of the budget and the lattice-theoretic
\cref{prop:tadpole:even}. The relevant files are \lean{DualScaleStream2/Flux/Tadpole.lean} and
\lean{DualScaleStream2/Flux/Integrality.lean}, and four files in older libraries that model
the orientifold charges by integer literals. All audited theorems depend only on the axioms
\lean{propext}, \lean{Classical.choice} and \lean{Quot.sound}.

\begin{leanbox}[In Lean: Euler numbers and the anomaly of $K3\times K3$]
\begin{leancode}
def eulerFromHodge {m : ℕ} (h : Fin m → Fin m → ℕ) : ℤ :=
  ∑ p : Fin m, ∑ q : Fin m,
    (if (p.val + q.val) % 2 = 0 then 1 else -1 : ℤ) * (h p q : ℤ)

theorem euler_K3 :
    eulerFromHodge StringTheory.Frontier.k3HodgeNumber = 24 := by ...

def chiK3K3 : ℤ := eulerFromHodge StringTheory.Frontier.k3HodgeNumber *
  eulerFromHodge StringTheory.Frontier.k3HodgeNumber

theorem k3k3_anomaly : chiK3K3 % 24 = 0 ∧ chiK3K3 / 24 = 24 := by ...
\end{leancode}
\emph{Reading.} \lean{eulerFromHodge} is the alternating sum $\sum(-1)^{p+q}h^{p,q}$ of an
arbitrary square table of natural numbers. \lean{k3HodgeNumber} is the $3\times3$ table
$(1,0,1;0,20,0;1,0,1)$, entered by hand in another library; \lean{euler\_K3} says its
alternating sum is $24$. The Euler number of $K3\times K3$ is \emph{defined} as the product
$24\cdot24$: multiplicativity (K\"unneth) is an input, not a theorem. \lean{k3k3\_anomaly}
then states that this integer is divisible by $24$ with quotient $24$. It does not know what
a K3 surface is, and it says nothing about $X_8$. \emph{Proof idea.} Unfold the double sum
over \lean{Fin 3}, evaluate, and finish with \lean{norm\_num} after rewriting with
\lean{euler\_K3}. The companions \lean{euler\_T2} and \lean{k3t2\_euler\_zero} are discussed
in \cref{ch:k3t2}.
\end{leanbox}

\begin{leanbox}[In Lean: the budget \eqref{eq:tadpole:budget}]
\begin{leancode}
theorem tadpole_budget (flux n : ℕ) (h : flux + n = 24) :
    n ≤ 24 ∧ (flux = 0 → n = 24) := by ...
\end{leancode}
\emph{Reading.} For natural numbers \lean{flux} and \lean{n} adding up to $24$, $n\le24$, and
$n=24$ if the flux term vanishes. The two physical inputs of \eqref{eq:tadpole:budget} are
encoded in the \emph{types}: \lean{flux : ℕ} is the positivity \eqref{eq:tadpole:pos} together
with the integrality of \cref{prop:tadpole:even}, and \lean{n : ℕ} is the absence of
anti-branes. The hypothesis \lean{h} is the tadpole condition itself, assumed. The theorem
constructs no flux and proves the existence of no vacuum. \emph{Proof idea.} Linear
arithmetic over $\N$, closed by \lean{omega}.
\end{leanbox}

\begin{leanbox}[In Lean: \cref{prop:tadpole:even}]
\begin{leancode}
def kronForm (L₁ : Gram n) (L₂ : Gram m) :
    Matrix (Fin n × Fin m) (Fin n × Fin m) ℤ :=
  kroneckerMap (· * ·) L₁ L₂

theorem kronForm_even (L₁ : Gram n) (L₂ : Gram m)
    (h₁s : L₁ᵀ = L₁) (h₂s : L₂ᵀ = L₂)
    (h₁ : IsEvenDiag L₁) (x : Fin n × Fin m → ℤ) :
    Even (x ⬝ᵥ (kronForm L₁ L₂ *ᵥ x)) := by ...

theorem flux_half_selfIntersection_integral (L₁ : Gram n) (L₂ : Gram m)
    (h₁s : L₁ᵀ = L₁) (h₂s : L₂ᵀ = L₂) (h₁ : IsEvenDiag L₁)
    (x : Fin n × Fin m → ℤ) :
    ∃ k : ℤ, x ⬝ᵥ (kronForm L₁ L₂ *ᵥ x) = 2 * k := by ...
\end{leancode}
\emph{Reading.} \lean{Gram n} is the type of $n\times n$ integer matrices and
\lean{IsEvenDiag L} says every diagonal entry is even. \lean{kronForm} is the matrix
\eqref{eq:tadpole:kron} indexed by pairs. The theorems say: for symmetric $L_1,L_2$ with $L_1$
even on the diagonal, and \emph{every} integer vector $x$ indexed by pairs,
$x^{\mathsf T}(L_1\otimes L_2)x$ is even, i.e.\ equals $2k$ for an integer $k$. This is a
theorem about all pairs of integer matrices of all sizes; it is genuinely general, and the
example \eqref{eq:tadpole:drsgram} is an instance. It is \emph{not} a statement about the
$22\times22$ K3 matrix specifically: to apply it to $3U\oplus2E_8(-1)$ one needs that Gram
matrix to be symmetric with even diagonal, for which the library provides the blocks
\lean{hyperbolicU\_evenDiag} and \lean{cartanE8\_evenDiag} (\cref{ch:lattices}). Nor does it
cover the rank-two summand of \eqref{eq:tadpole:H4}, or say that fluxes are elements of this
lattice; those are Tier~L. \emph{Proof idea.} The pair index is flattened by the bijection
\lean{finProdFinEquiv} between \lean{Fin n × Fin m} and \lean{Fin (n * m)}; symmetry and even
diagonal survive the relabelling (\lean{kronForm\_symm}, \lean{kronForm\_evenDiag}); then the
single-index lemma \lean{even\_quadratic\_form\_of\_even\_diag}, which is the computation in
the proof of \cref{prop:tadpole:even}, is applied and the result transported back.
\end{leanbox}

\begin{leanbox}[In Lean: orientifold charges in the older libraries (arithmetic shadows)]
\begin{leancode}
-- StringTheory.StringDynamics (TadpoleConstraint.lean)
def totalTadpole (braneStacks : Fin 4 → BraneStack) (H₃ F₃ : ℤ) : ℤ :=
  (∑ i, (braneStacks i).charge) + fluxTadpole H₃ F₃ - 24
theorem tadpole_cancellation (braneStacks : Fin 4 → BraneStack) (H₃ F₃ : ℤ)
    (h : totalTadpole braneStacks H₃ F₃ = 0) :
    (∑ i, (braneStacks i).charge) + fluxTadpole H₃ F₃ = 24 := by ...

-- StringTheory.Foundation.StringTheory.TadpoleCancellation
theorem d7_tadpole_cancellation : totalD7Charge + totalO7Charge = 0 := by rfl
theorem d3_tadpole_target_is_one : d3TadpoleTarget = 1 := by rfl

-- SocrateAI.Moonshine (KummerTadpole.lean)
theorem rr_tadpole_cancellation : net_tadpole_charge = 0 := by rfl
-- DualScaleValidation.UseCase3, with its own copies of the definitions
theorem rr_tadpole_cancellation :
    total_d7_charge + total_o7_charge = 0 := rfl
\end{leancode}
\emph{Reading.} These are statements about integer literals. \lean{fluxTadpole} is the
product of two integers; \lean{tadpole\_cancellation} says that if
$\sum_i q_i+H_3F_3-24=0$ then $\sum_iq_i+H_3F_3=24$, a rearrangement of its own hypothesis
(the \lean{euler} field of \lean{BraneStack} is never used).
\lean{d7\_tadpole\_cancellation} unfolds to $32\cdot2+16\cdot(-4)=0$ and
both \lean{rr\_tadpole\_cancellation} to $16\cdot4+4\cdot(-16)=0$; all are proved by
\lean{rfl}.
Compare with \eqref{eq:tadpole:d7}: the second is the $p=7$ row of
\cref{tab:tadpole:oplanes} ($4$ planes, $16$ branes, ratio $-4$) with all charges multiplied
by $4$. The first has $16$ planes; by \eqref{eq:tadpole:oplane} sixteen planes belong to
$p=5$, where the ratio of plane to brane charge is $-1$, and $32\cdot2$ against $16\cdot(-4)$
is that row with branes and images counted separately and charges again multiplied by $4$, so
the label ``O7'' in that file does not match its numbers. This is the convention issue of the
pitfall box in \cref{sec:tadpole:gauss}; the Lean statements themselves are correct
arithmetic and carry no information about which convention is meant.
\lean{d3\_tadpole\_target\_is\_one} is $24/24=1$, the number that appears in
\cref{rem:tadpole:where}. Finally,
\begin{leancode}
theorem flux_tadpole_quantization_exact
    (s : StringTheory.Foundation.StringTheory.VafaSwampland.GVWFluxState) :
    StringTheory.Foundation.StringTheory.TadpoleCancellation.totalD7Charge +
    StringTheory.Foundation.StringTheory.TadpoleCancellation.totalO7Charge
      = 0 := ...
\end{leancode}
(in \lean{StringTheoryFormalization/Frontier/FTermPotential.lean}) never uses its argument
\lean{s} and is proved by quoting \lean{d7\_tadpole\_cancellation}. Despite its name it
contains no flux and no quantization. A faithful formalization of \eqref{eq:tadpole:d7} would
need at least the fixed-point set of the involution on a torus as a finite type of
cardinality $2^{9-p}$ and the charge ratio $-2^{p-5}$ as a function of $p$, so that the
independence of the total from $p$ becomes a theorem; see \cref{exe:tadpole:leanoplane}.
\end{leanbox}

\paragraph{Unification candidates.} The atlas of \cref{ch:atlas} compares statements across
libraries by text similarity and by overlap of their dependencies. Under the heading
``similar statements, essentially disjoint proofs'' it lists two pairs from this chapter
\source{\texttt{papers/book/generated/atlas.md}, ll.~254--262}:
\begin{itemize}[leftmargin=2em,nosep]
\item \lean{rr\_tadpole\_cancellation} (in \lean{UseCase3}) and
 \lean{tadpole\_cancellation}\\ (in \lean{StringDynamics}): similarity $0.514$,
 dependency overlap $0.021$;
\item \lean{k3t2\_euler\_zero} (in \lean{Flux}) and \lean{k3t2\_euler\_char\_eq\_zero}
 (in \lean{Moonshine}):\\ similarity $0.557$, dependency overlap $0.009$.
\end{itemize}
The reading
above explains the finding: the tadpole of this chapter has been written down four times with
four sets of private definitions. A single definition of the budget, with the Euler number
supplied by \lean{eulerFromHodge} and the flux term by \lean{kronForm}, would let these
statements share content.

{\small
\begin{longtable}{@{}p{0.44\textwidth}cp{0.43\textwidth}@{}}
\toprule
Claim & Tier & Evidence\\
\midrule
Alternating sum of the K3 Hodge table is $24$; $24^2\bmod24=0$, $24^2/24=24$ & \TA &
 \lean{euler\_K3}, \lean{k3k3\_anomaly}\\
$\mathit{flux}+n=24$ in $\N$ $\Rightarrow$ $n\le24$, and $n=24$ if $\mathit{flux}=0$ & \TA &
 \lean{tadpole\_budget}\\
$L_1$ even-diagonal, $L_1,L_2$ symmetric $\Rightarrow$ $x^{\mathsf T}(L_1\otimes L_2)x$ even
 for all $x$ & \TA & \lean{kronForm\_even},\newline
 \lean{flux\_half\_selfIntersection\_integral}\\
$64-64=0$ in two bookkeepings; $24/24=1$ & \TA & \lean{rr\_tadpole\_cancellation},
 \lean{d7\_tadpole\_cancellation}, \lean{d3\_tadpole\_target\_is\_one}\\
Betti and Hodge numbers of $K3\times K3$, \eqref{eq:tadpole:div6}, $2610$, $972$, $-192$,
 signature $(371,115)$, table of \cref{sec:tadpole:worked} & --- & symbolic checks (sympy), not
 Tier~A\\
$I=-\int X_8=\chi/24$; tadpole condition \eqref{eq:tadpole:master}; obstructions & \TL &
 SVW §1; BB §2; DRS §1--2; GVW §2.1\\
$6\mid\chi$; \eqref{eq:tadpole:elliptic}; shifted quantization of $G$; supersymmetric $G$
 primitive and self-dual & \TL & SVW §2.1--2.2; GVW §2.1, §5; DRS §2.1\\
O$p$-plane charge $-2^{p-5}\mu_p$; $SO(8)^4$; $24$ D3 $\leftrightarrow$ $24$ D5
 $\leftrightarrow$ $24$ five-branes & \TL & Pol Lecture~3; DRS §2.4, §4; SVW §2.4\\
``\lean{kronForm L L} \emph{is} the flux lattice of M-theory on $K3\times K3$'' & --- &
 an identification, never Tier~A\\
\bottomrule
\end{longtable}}

\begin{summarybox}
\begin{itemize}[leftmargin=1.5em]
\item On a compact space Gauss's law forces the total charge of every gauge field to vanish.
 A violation appears in perturbation theory as a one-point function, a tadpole, and means
 that the field equations have no solution.
\item In compactifications on a Calabi--Yau fourfold the curvature carries charge
 $\chi/24$ through the interaction $C\wedge X_8$. It must be balanced by $n$ space-filling
 branes and by flux: $\chi/24=n+\tfrac12\int\frac G{2\pi}\wedge\frac G{2\pi}$.
\item Supersymmetry requires $n\ge0$ and a primitive, hence self-dual, $G$, so both terms are
 non-negative: a budget, equal to $576/24=24$ for $K3\times K3$. Branes alone suffice only if
 $\chi/24$ is a non-negative integer;
 $\chi$ is always divisible by $6$ but not always by $24$ (sextic: $108\tfrac34$), and can be
 negative in orbifolds with discrete torsion ($-8$).
\item Flux is quantized, $[G/2\pi]\in H^4(N;\Z)$ up to a shift by $c_2/2$ that is absent when
 the intersection form is even. On $H^2(K3)\otimes H^2(K3)$ the form is the Kronecker product
 $L\otimes L$, which is even because $L$ is; hence $\tfrac12\int\xi\wedge\xi\in\Z$.
\item In type IIB on $\KT/\Z_2$ there are two tadpoles. The seven-brane tadpole is cancelled
 locally, $4\times(-4)+16=0$, giving $SO(8)^4$. The three-brane tadpole is $24$ and is shared
 between D3-branes and $H\wedge H'$ flux; without flux it is T-dual to type~I with $24$
 D5-branes.
\item Lean checks the arithmetic ($576/24=24$, the budget inequality, $64-64=0$) and the
 general evenness of Kronecker forms. The physical identifications are Tier~L.
\end{itemize}
\end{summarybox}

\section*{Exercises}
\addcontentsline{toc}{section}{Exercises}

\begin{exercise}[$\star$]
Two point charges $q_1,q_2$ sit on a circle of circumference $\ell$ and the electric field
obeys $\partial_xE=\sum_iq_i\delta(x-x_i)$. Show directly, by solving for the piecewise
constant $E(x)$ and demanding periodicity, that a solution exists if and only if
$q_1+q_2=0$. Which step fails on the real line?
\end{exercise}

\begin{exercise}[$\star$]
Compute the Betti numbers and the Euler number of $K3\times T^4$ and of $T^8$ from Poincar\'e
polynomials. What does the tadpole condition say about M-theory on these spaces with
unbroken supersymmetry? (Use \eqref{eq:tadpole:budget}.)
\end{exercise}

\begin{exercise}[$\star\star$]\label{exe:tadpole:hodge}
A Calabi--Yau fourfold with holonomy exactly $SU(4)$ has independent Hodge numbers
$h^{1,1},h^{2,1},h^{3,1}$ and $h^{2,2}$. (a) Write $\chi_1$ and $\chi$ in terms of them.
(b) Use $\chi_1=8-\chi/6$ to derive $h^{2,2}=2\,(22+2h^{1,1}+2h^{3,1}-h^{2,1})$ and
$\chi=6\,(8+h^{1,1}+h^{3,1}-h^{2,1})$. (c) The sextic has $h^{1,1}=1$, $h^{2,1}=0$,
$h^{3,1}=426$. Recover $\chi=2610$.
\end{exercise}

\begin{exercise}[$\star\star$]
For the lattice \eqref{eq:tadpole:drsgram}, compute $\tfrac12\int\xi\wedge\xi$ for the nine
decomposable fluxes $\xi=a_1\wedge a_2$ with $a_1,a_2\in\{\alpha,\beta,\alpha+\beta\}$. Which
of them fit in the budget, and with how many membranes? Show that every integral class
$a=xC_1+yC_2+zC_3$ orthogonal to $J$ has $a\cdot a=-2(x^2+y^2+z^2)$, and conclude that only
finitely many decomposable primitive fluxes satisfy \eqref{eq:tadpole:budget}.
\end{exercise}

\begin{exercise}[$\star\star$]
Show that the rank-two summand in \eqref{eq:tadpole:H4} has Gram matrix
$\bigl(\begin{smallmatrix}0&1\\1&0\end{smallmatrix}\bigr)$ and is orthogonal to
$\Lambda_1\otimes\Lambda_2$. Deduce the signature \eqref{eq:tadpole:sig} and check it against
$\sigma(X\times Y)=\sigma(X)\sigma(Y)$. Is $U\oplus(\Lambda\otimes\Lambda)$ unimodular?
(Hint: $\det(A\otimes B)=(\det A)^{m}(\det B)^{n}$.)
\end{exercise}

\begin{exercise}[$\star\star$, Lean]
State and prove in Lean the strengthening of \lean{tadpole\_budget} that is symmetric in its
two arguments: for \lean{flux n : ℕ} with \lean{flux + n = 24}, also \lean{flux ≤ 24} and
\lean{n = 0 ↔ flux = 24}. Which tactic closes it? Then explain why the statement would be
false for \lean{flux n : ℤ} and what physical assumption the type \lean{ℕ} encodes.
\end{exercise}

\begin{exercise}[$\star\star\star$, Lean]
Instantiate \lean{flux\_half\_selfIntersection\_integral} with
$L_1=L_2=L_{\mathrm{DRS}}$ of \eqref{eq:tadpole:drsgram}, written as
\lean{!![0,2,2;2,0,2;2,2,0]}: prove the two symmetry hypotheses and \lean{IsEvenDiag} (by
\lean{decide} or \lean{fin\_cases}). Then state, for \lean{a b : Fin 3 → ℤ} and
\lean{x := fun p => a p.1 * b p.2}, the factorization
$x^{\mathsf T}(L_1\otimes L_2)x=(a^{\mathsf T}L_1a)(b^{\mathsf T}L_2b)$ and prove it for
general \lean{Gram n} and \lean{Gram m}. (Hint: use \lean{Matrix.kroneckerMap\_apply}
and\\ \lean{Finset.sum\_mul\_sum}.)
\end{exercise}

\begin{exercise}[$\star\star\star$, Lean]\label{exe:tadpole:leanoplane}
Replace the integer literals of \lean{d7\_tadpole\_cancellation} by a statement with content.
Work in units of a quarter of $\mu_p$ so that all charges are integers for $3\le p\le9$:
define \lean{planes p := 2 \^{} (9 - p)} and \lean{planeCharge p := -(2 \^{} (p - 3) : ℤ)} and prove
that \lean{(planes p : ℤ) * planeCharge p + 16 * 4 = 0} for every \lean{p} with
\lean{3 ≤ p} and \lean{p ≤ 9}. Identify which of the two older bookkeepings
($16\cdot4-4\cdot16$ and $32\cdot2-16\cdot4$) your statement reproduces at $p=7$ and at $p=5$.
\end{exercise}

\section*{Further reading}
\addcontentsline{toc}{section}{Further reading}
\begin{itemize}[leftmargin=1.5em]
\item The tadpole $\chi/24$, its two obstructions, divisibility by six, the elliptic formula
 and the $K3\times K3$ duality check: \source{SVW §1--2, ll.~34--546}. Its Section~3 notes
 that blowing up a point of the base removes $120$ branes \source{SVW §3, ll.~571--588}.
\item The supergravity solution with warp factor and the flux constraint from the Bianchi
 identity: \source{BB §2, eqs.~(2.6)--(2.7), (2.56)--(2.59), ll.~117--146, 668--735}.
\item $G$-flux on $K3\times K3$, the lift to F-theory, the orientifold of $\KT$ and its type~I
 and heterotic duals: \source{DRS §1--4, ll.~53--148, 156--266, 457--578, 636--808,
 1262--1400}.
\item Flux quantization, its shift, vacua as states of one model, and the F-theory analogue
 of $G$: \source{GVW §2.1, ll.~84--135; §2.5, ll.~1068--1111}. The superpotential generated by
 $G$ is the subject of \cref{ch:stabilization}.
\item Orientifold-plane charges and the sixteen D-branes: \source{Pol Lecture~3,
 ll.~1968--2000}. D3 charge induced on wrapped seven-branes and the integrated Bianchi
 identity in type IIB: \source{GKP §2.1, ll.~419--502}.
\end{itemize}
